\section{Methodology: Discrete Entropy Lower Bound}
\label{sec:theory}

\subsection{General Integer-Valued Prediction Setting}
\label{subsec:notation}

Let \(\{X_t\}_{t\in\mathbb{Z}}\) be an integer-valued stochastic process,
\begin{equation}
X_t\in\mathbb{Z}.
\label{eq:general_integer_process}
\end{equation}
For an \(m\)-step forecast, \(m\geq1\), let
\(\mathcal{F}_{t-m}\) denote the sigma-algebra containing exactly the
observations available when the forecast is issued. A non-anticipating
integer-valued predictor is selected from the explicitly restricted decision class
\begin{equation}
\mathcal A_{\mathbb Z}(\mathcal F_{t-m})
=
\left\{
f(\mathcal F_{t-m}):
f\ \text{is measurable},\
f(\mathcal F_{t-m})\in\mathbb Z,\
\mathbb E[(X_t-f(\mathcal F_{t-m}))^2]<\infty
\right\},
\label{eq:general_predictor}
\end{equation}
and \(\widehat X_t\in\mathcal A_{\mathbb Z}(\mathcal F_{t-m})\).
Its prediction error and mean squared error (MSE) are
\begin{equation}
E_t=X_t-\widehat X_t,
\qquad
D_t^{(m)}=\mathbb{E}[E_t^2].
\label{eq:general_error_mse}
\end{equation}
\rev{Here, \(E_t\) is the integer-valued forecast error, \(X_t\) is the
realized count, \(\widehat X_t\) is the forecast issued from information
available at time \(t-m\), and \(D_t^{(m)}\) is the expected squared forecast
error at horizon \(m\). Thus, every MSE bound below applies to the prediction
error \(X_t-\widehat X_t\), not to the count itself or to a parameter
estimator.}
All logarithms have base \(2\); entropy is measured in bits.

\begin{Assumption}
\label{ass:general_regularity}
For every time and horizon considered,
\(H(X_t\mid\mathcal{F}_{t-m})<\infty\) and
\(\mathbb{E}[E_t^2]<\infty\). Forecast information sets are temporally admissible and
nested:
\(\mathcal{F}_{t-m}\subseteq\mathcal{F}_{t-\ell}\) whenever
\(1\leq\ell<m\).
\end{Assumption}

For a discrete random variable \(V\) with mass function \(p(v)\),
\begin{equation}
H(V)=-\sum_vp(v)\log_2p(v),
\label{eq:shannon_entropy}
\end{equation}
\rev{Here, \(p(v)=\Pr(V=v)\) is the probability mass at outcome \(v\).}
and conditional mutual information (CMI) satisfies
\begin{equation}
I(V;W\mid Y)=H(V\mid Y)-H(V\mid Y,W).
\label{eq:cmi_definition}
\end{equation}
\rev{Here, \(V\) is the target, \(W\) is the added information, and \(Y\) is
the information already available for conditioning.}
The distinction between a population quantity such as
\(H(X_t\mid\mathcal F_{t-m})\) and its finite-sample estimate is essential:
the former defines the lower bound below, whereas the latter is an inferential
object studied in Sections~\ref{sec:methods} and \ref{sec:results}.

\subsection{Discrete Entropy Prediction-Error Lower Bound and Predictor-Dependent Refinement}
\label{subsec:discrete_bound}

Continuous entropy--variance arguments rely on the fact that a Gaussian
maximizes differential entropy \citep{fang_variance_bounds}. An
integer-valued error lies on a lattice, so differential entropy cannot be
replaced mechanically by Shannon entropy. The relevant extremal object is
the integer-lattice maximum-entropy envelope
\citep{massey_integer_entropy}.

\begin{Definition}[Integer maximum-entropy envelope]
\label{def:integer_entropy_envelope}
For \(D\geq0\), define
\begin{equation}
\mathcal H_{\mathbb Z}(D)
=
\sup_{\substack{P_Z:\,Z\in\mathbb Z\\\mathbb E[Z^2]\leq D}}H(Z),
\label{eq:integer_entropy_envelope}
\end{equation}
\rev{Here, \(P_Z\) ranges over distributions of integer-valued random
variables \(Z\), and \(D\) is an upper bound on their second moment.}
with generalized inverse
\begin{equation}
\mathcal H_{\mathbb Z}^{-1}(h)
=
\inf\{D\geq0:\mathcal H_{\mathbb Z}(D)\geq h\}.
\label{eq:integer_entropy_inverse}
\end{equation}
\end{Definition}

For \(D>0\), the envelope is attained by the centered discrete Gaussian
\begin{equation}
q_\lambda(k)
=
\frac{\exp(-\lambda k^2)}{\Theta(\lambda)},
\qquad
\Theta(\lambda)=\sum_{j\in\mathbb Z}\exp(-\lambda j^2),
\quad k\in\mathbb Z,
\label{eq:general_discrete_gaussian}
\end{equation}
\rev{Here, \(k\) is an integer error value, \(\lambda>0\) controls
dispersion, and \(\Theta(\lambda)\) normalizes the mass function.}
where \(\lambda>0\) is selected so that
\begin{equation}
D=D(\lambda)
=-\frac{\mathrm d}{\mathrm d\lambda}\ln\Theta(\lambda).
\label{eq:general_discrete_gaussian_moment}
\end{equation}
Its entropy is
\begin{equation}
H_\lambda
=
\frac{\lambda D(\lambda)+\ln\Theta(\lambda)}{\ln2}.
\label{eq:general_discrete_gaussian_entropy}
\end{equation}
Equations~\eqref{eq:general_discrete_gaussian_moment} and
\eqref{eq:general_discrete_gaussian_entropy} define the one-dimensional
numerical inversion used in this study.

\begin{Theorem}[Discrete entropy prediction-error lower bound and predictor-dependent refinement]
\label{thm:exact_delb}
Under Assumption~\ref{ass:general_regularity}, every non-anticipating integer-valued
predictor satisfies
\begin{align}
D_t^{(m)}
&\geq
\mathcal H_{\mathbb Z}^{-1}\!\left[
H(X_t\mid\mathcal F_{t-m})
+I(E_t;\mathcal F_{t-m})
\right]
\label{eq:general_predictor_dependent_delb}\\
&\geq
\mathcal H_{\mathbb Z}^{-1}\!\left[
H(X_t\mid\mathcal F_{t-m})
\right].
\label{eq:general_universal_delb}
\end{align}
The second line is the discrete entropy prediction-error lower bound
(DELB).
\end{Theorem}

Equation~\eqref{eq:general_universal_delb} gives the population-level answer
to RQ1: target uncertainty and the integer lattice impose an MSE lower bound
independently of the forecasting algorithm.
Equation~\eqref{eq:general_predictor_dependent_delb} is stronger for a
particular predictor because it retains the residual information quantified
by \(I(E_t;\mathcal F_{t-m})\). The proof and its boundary cases accompany
the article together with the closed-form, temporal, stationary, and vector
consequences.

To distinguish the DELB from minimum decision risk, define the optimal risk
within this same non-anticipating integer-valued predictor class by
\begin{equation}
R_{\mathbb Z}^{*}(\mathcal F)
=
\inf_{\widehat X\in\mathcal A_{\mathbb Z}(\mathcal F)}
\mathbb E[(X-\widehat X)^2].
\label{eq:integer_optimal_risk}
\end{equation}
For every admissible integer-valued predictor,
\begin{equation}
L_{\mathrm{DELB}}(\mathcal F)
\leq
R_{\mathbb Z}^{*}(\mathcal F)
\leq
\mathbb E[(X-\widehat X)^2].
\label{eq:delb_bayes_risk_hierarchy}
\end{equation}
Thus, the exactness of the DELB concerns the discrete entropy envelope and
its inversion. It does not make the DELB identical to optimal risk. If the
infimum in Equation~\eqref{eq:integer_optimal_risk} is attained, equality
\(L_{\mathrm{DELB}}=R_{\mathbb Z}^{*}\) holds if and only if a minimizing
predictor satisfies the attainability conditions below. Without attainment
of the infimum, the hierarchy remains valid, but equality of numerical
infima must not be interpreted as the existence of a predictor that attains
the DELB.

\subsection{Attainability and the Two-Slack Decomposition}
\label{subsec:attainability}

\begin{Proposition}[Exact entropy-domain attainability decomposition]
\label{prop:attainability_decomposition}
Abbreviate \(\mathcal F=\mathcal F_{t-m}\), \(E=E_t\), and
\(D=\mathbb E[E^2]\). For \(D>0\), let \(q_D\) be the centered discrete
Gaussian with second moment \(D\); for \(D=0\), let \(q_0\) be the point mass
at zero. Then
\begin{align}
\mathcal H_{\mathbb Z}(D)
-\left[H(X_t\mid\mathcal F)+I(E;\mathcal F)\right]
&=
D_{\mathrm{KL},2}(P_E\Vert q_D),
\label{eq:predictor_attainability_decomposition}\\
\mathcal H_{\mathbb Z}(D)-H(X_t\mid\mathcal F)
&=
\underbrace{I(E;\mathcal F)}_{\text{history-dependence slack}}
+
\underbrace{D_{\mathrm{KL},2}(P_E\Vert q_D)}_
{\text{shape-mismatch slack}},
\label{eq:universal_attainability_decomposition}
\end{align}
where \(D_{\mathrm{KL},2}\) is the Kullback--Leibler (KL) divergence in bits.
\end{Proposition}

The predictor-specific bound is attained if and only if \(P_E=q_D\).
The universal DELB is attained if and only if \(P_E=q_D\) and
\begin{equation}
I(E_t;\mathcal F_{t-m})=0.
\label{eq:general_innovation_independence}
\end{equation}
The error must be moment-matched discrete Gaussian and independent of the
admissible history. Zero correlation or second-order residual
whiteness is insufficient. The two terms in
Equation~\eqref{eq:universal_attainability_decomposition} diagnose different
population sources of nonattainment; neither is the empirical
Predictability Gap introduced in Section~\ref{subsec:methods_prediction}.

Detailed comparisons with neighboring information-theoretic bounds and the
secondary corollaries are kept outside the main text so that it can focus on
the implementable bound, its attainability diagnostics, and added-information
ordering.

\begin{Remark}[Boundary and loss conditions]
\label{rem:general_boundary_cases}
Because \(\mathcal H_{\mathbb Z}^{-1}(0)=0\), information that determines
\(X_t\) almost surely produces a zero DELB. Conditionally irrelevant side
information leaves it unchanged. Shannon entropy alone does not define an
MSE lower bound without integer support and squared-error loss; the lattice
envelope is an essential component of Theorem~\ref{thm:exact_delb}.
\label{rem:support_and_loss}
\end{Remark}

\subsection{Added Side Information and CMI Ordering}
\label{subsec:side_information}

Let \(\mathcal F_t^0\) be a baseline forecast-time information set and let \(S_t\)
be additional side information observed when the forecast is issued. Define
\begin{equation}
\mathcal F_t^S=\mathcal F_t^0\vee\sigma(S_t).
\label{eq:generic_augmented_information}
\end{equation}
\rev{Here, \(\mathcal F_t^0\) is the baseline forecast-time information,
\(S_t\) is the added observation, \(\sigma(S_t)\) is the sigma-algebra it
generates, and \(\vee\) denotes the information generated jointly by its
arguments.}
The population uncertainty reduction is exactly
\begin{equation}
\Delta H^S_t
=
H(X_t\mid\mathcal F_t^0)-H(X_t\mid\mathcal F_t^S)
=
I(X_t;S_t\mid\mathcal F_t^0)\geq0.
\label{eq:generic_information_gain}
\end{equation}
Because the inverse envelope is nondecreasing, the corresponding universal floors
\begin{equation}
L_t^0=\mathcal H_{\mathbb Z}^{-1}
\!\left[H(X_t\mid\mathcal F_t^0)\right],
\qquad
L_t^S=\mathcal H_{\mathbb Z}^{-1}
\!\left[H(X_t\mid\mathcal F_t^S)\right]
\label{eq:generic_side_information_bounds}
\end{equation}
satisfy
\begin{equation}
L_t^S\leq L_t^0.
\label{eq:generic_side_information_ordering}
\end{equation}

Equations~\eqref{eq:generic_information_gain}--
\eqref{eq:generic_side_information_ordering} define the population estimand
for RQ2 and RQ3. They do not imply that an estimated CMI is
unbiased, that a design-conditioned randomization distribution is centered
at zero, or that a finite-sample predictor will realize the population
reduction. These are separate empirical questions. Deterministic
coarsening \(Q(S_t)\) cannot increase conditional information:
\begin{equation}
I(X_t;Q(S_t)\mid\mathcal F_t^0)
\leq
I(X_t;S_t\mid\mathcal F_t^0).
\label{eq:generic_data_processing}
\end{equation}

The Supplementary Materials give the complete proofs, place the DELB relative
to established information-theoretic bounds, and state the supplementary
classification-loss result. Section~\ref{sec:methods} maps the generic objects
\((X_t,\mathcal F_t^0,S_t)\) to crime counts, admissible controls, and lagged
bicycle mobility.

The four main-text components have distinct operational roles: the universal
DELB defines the benchmark, the predictor-dependent term and two-slack
decomposition diagnose nonattainment, and CMI orders information sets at the
population level. Section~\ref{sec:methods} supplies the finite-sample
estimation, calibration, and matched-prediction workflow needed to apply
these quantities.
